\chapter{Questions to Astma-Allergi Danmark}\label{ch:experts}
% Draft by Claude, edit freely. Everything in this chapter is a plan. Nothing has been done yet.
% Rewrite it in past tense with the answers once the email is sent and answered.

\fillin{This chapter is a plan. The email has not been sent yet. Update it with what we
actually did and found.}

Astma-Allergi Danmark runs the pollen counts in Denmark today~\cite{AstmaAllergi}. If cheap
stations are ever going to help the forecasts, the data has to be something they can use and
trust. Before we design more of the device, we want to know what they need from it.

\section{Written questions}\label{sec:expert-email}

\paragraph{Goal.} Learn how the people who count pollen today work, and what they would need
before they could use data from cheap stations.

\paragraph{Method.} A meeting is hard to arrange, so we send a short email with written
questions (Appendix~\ref{app:expert-questions}). This is a written interview: the answers are
shorter than in a conversation and we can only ask follow-up questions in a new email, but it
takes little of their time, which makes an answer more likely.
\fillin{Who we sent it to, when, and whether they answered.}

\paragraph{Topics.}
\begin{itemize}
  \item How the daily count is made, and how much work it takes.
  \item What accuracy and time resolution a count needs to be useful for a forecast.
  \item What they would need to see before they trust an automated count.
  \item Whether a network of cheap stations would be useful to them, and who would look after
  the stations.
\end{itemize}

\section{What we learned}\label{sec:expert-results}

\fillin{The answers, and what we change in the design because of them. Write this when the
answers are in.}
